\section{A localized explicit formula}
\label{sec:localized-kernel}

Fix $b>1/\pi$ and $1/2<d<\pi b/2$ (one may fix the midpoint).
Use $\sinc v=\sin(v)/v$ with value $1$ at zero. Put
\[
 k(t)=e^{-\cosh(t/b)},\qquad
 \eta(v)=\frac{\int_{\R}k(t)e^{-itv}\,dt}{\int_{\R}k(t)\,dt},
 \qquad K_B(v)=\eta(v)\sinc^4(Bv).
\]
For a prime $p$, write $\ell=\log p$. For $B\ge1$ and $U\in\R$ define
\[
 g(u)=\frac{u}{\ell}e^{-iU(u-\ell)}K_B(u-\ell),\quad
 h(t)=\int_{\R}g(u)e^{itu}\,du,\quad
 D_{p,B,U,H}(\chi)=-\frac{\sqrt p}{\ell}
       \sum_{|\gamma-U|\le H}h(\gamma).
\]
The ordinates are signed and counted with multiplicity. All constants below
may depend on $b,d$. Write $\delta_\chi=1$ for the conductor-one character and
$0$ otherwise. Define the nonnegative leakage from nearby prime powers by
\[
 L_B(p)=\sum_{\substack{p/e\le n\le ep,\ n\ne p\\ \Lambda(n)>0}}
   \left(1+B\frac{|n-p|}{p}\right)^{-4}.
\]

\begin{lemma}[Prime-value decoder]\label{lem:kernel-decoder}
For primitive $\chi$ of conductor $q\le Q$, under Dirichlet GRH, and for
$H=4B+\Delta$, $\Delta\ge1$, one has
\begin{align*}
 |D_{p,B,U,H}(\chi)-\chi(p)|\ll_{b,d}{}&L_B(p)+\frac{p}{B^4}+\frac1B
       +\delta_\chi\frac{p}{B\ell}\\[-2pt]
 &+\frac{\sqrt p}{\ell}\log\!\big(eQ(3+|U|+H)\big)
                        \exp\{-c_b e^{\Delta/b}\}
\end{align*}
with some $c_b>0$. Given $\epsilon>0$, the last term, with its implied
constant, is $\le\epsilon$ whenever
\begin{equation}\label{eq:kernel-tail-choice}
 \Delta\ge C_{b,d,\epsilon}+b\log\!\big(2+\log(2p)+
           \log\log(e^e Q(3+|U|+B+p))\big).
\end{equation}
The normalized summand $-\sqrt p\,h(t)/\ell$ is bounded by
$O_{b,d}(\sqrt p/(B\ell))$ and has Lipschitz constant
$O_{b,d}(\sqrt p/B)$ on $\R$.
\end{lemma}

\begin{proof}
The defining integral gives $|\eta(v)|\le1$. Shifting its contour by
$-id\operatorname{sgn} v$ yields $|\eta(v)|\ll_{b,d}e^{-d|v|}$, since
$|k(t\pm id)|=\exp(-\cos(d/b)\cosh(t/b))$. Thus $h$ is analytic beyond
the closed strip $|\operatorname{Im}t|\le1/2$. In our Fourier convention,
the Guinand--Weil formula \cite[Eqs.~(2.1) and (2.6)]{weil} reads, when $g(0)=0$,
\begin{align*}
 \sum_\gamma h(\gamma)={}&\delta_\chi\{h(i/2)+h(-i/2)\}\\
 &-\int_0^\infty \frac{e^{-(a+1/2)u}}{1-e^{-2u}}
                 \{g(u)+g(-u)\}\,du\\
 &-\sum_{n\ge2}\frac{\Lambda(n)}{\sqrt n}
            \{\chi(n)g(\log n)+\overline{\chi(n)}g(-\log n)\},
\end{align*}
where $a\in\{0,1\}$ is the parity. For the sign convention, Mellin inversion
of $\int g(u)e^{(s-1/2)u}\,du$ against $-L'/L(s,\chi)$ on the right of $1$
selects $g(\log n)$. On the reflected line the functional equation supplies
$\overline\chi$ and $g(-\log n)$. The conductor term is a multiple of
$g(0)$, and the integral follows from the digamma integral by substituting
$t=2u$. Admissibility follows since $k$ and all its derivatives decrease
doubly exponentially in each smaller horizontal strip
$|\operatorname{Im}t|<\pi b/2$, and the transform of $K_B$ is an average of
translates of $k$ over a compact interval, as shown below.

We have $g(\ell)=1$. In the positive-frequency sum with
$|\log(n/p)|\le1$, the quantities $\sqrt{p/n}$, $|\log n|/\ell$ and
$\Lambda(n)/\ell$ are bounded absolutely, and
$|\log(n/p)|\ge |n-p|/(ep)$. Hence the normalized error there is
$O(L_B(p))$. For
$\log(n/p)\in[j,j+1)$, $j\ge1$, there are $O(pe^j)$ integers and each
normalized term is
$O_{b,d}((1+j)^2 e^{-(d+1/2)j}/B^4)$. Their sum is $O_{b,d}(p/B^4)$.
For $n<p/e$, the ratios $(\log n)/\ell$ and $\Lambda(n)/\ell$ are at most
one and $\sqrt{p/n}\,e^{-d|\log(n/p)|}\le1$, giving the same bound.
At negative frequencies $|\log(np)|\ge\log4$, so the error is at most
\[
 \frac{C_{b,d}p^{1/2-d}}{B^4\ell^2}
       \sum_{n\ge2}\frac{(\log n)^2}{n^{1/2+d}}\ll_{b,d}B^{-4}.
\]
For either sign of $u$ the archimedean density multiplied by $|u|$ is
$O((1+|u|)e^{-|u|/2})$. Using
$e^{-|u|/2}\le p^{-1/2}e^{|u-\ell|/2}$ and the decay of $\eta$ bounds
the normalized integral by $O_{b,d}(\int_\R\sinc^4(Bv)\,dv)=O_{b,d}(1/B)$.
The same computation with $e^{|u|/2}\le\sqrt p\,e^{|u-\ell|/2}$ bounds
the normalized pole terms by $O_{b,d}(p/(B\ell))$. Integrating $|g(u)|$ and
$|ug(u)|$ similarly proves the asserted size and Lipschitz bounds.

For localization, put $F(t)=\int K_B(v)e^{itv}\,dv$. The characteristic
function of the sum $V$ of four independent uniform variables on $[-B,B]$
is $\sinc^4(Bv)$. Fourier inversion gives
\[
 F(t)=\frac{2\pi}{\int k}\,\mathbb E[k(t+V)],\qquad
 h(t)=e^{it\ell}\{F(t-U)+F'(t-U)/(i\ell)\}.
\]
Both $|k(v)|$ and $|k'(v)|$ are $\ll_b\exp(-c_b e^{|v|/b})$. In the two
unit intervals with $H+j<|t-U|\le H+j+1$, $j=0,1,\ldots$, the standard
local zero count is $O(\log(eQ(3+|U|+H+j)))$. Throughout these intervals
$|t-U+V|\ge \Delta+j$. Summation in $j$ gives the stated tail estimate (decreasing
$c_b$ if necessary). To see \eqref{eq:kernel-tail-choice}, first take $\Delta$
at its displayed lower bound. Then $\log(eQ(3+|U|+H))$ is bounded by a
constant times $\log(e^eQ(3+|U|+B+p))$ and increasing $C_{b,d,\epsilon}$
absorbs this constant and $\sqrt p/\ell$. For larger $\Delta$ the omitted
absolute tail can only decrease.
\end{proof}
